\begin{frame}{5항목의 구체적 의미 (1) 환경 설명 · (2) 알고리즘 선택}
  \small
  \begin{itemize}
    \item \kb{(1) 환경 설명} — MDP 5요소(Ch03.1) 중 \textbf{상태·행동·보상}이
      보고서에서 추적 가능한지. ``Pendulum을 썼다'' 한 줄만 있고 \emph{주
      지표가 무엇인지}(학습 도중 리턴? 이동평균? 평가 리턴?), \emph{리턴의
      부호 규약}(항상 음수, 0에 가까울수록 좋음)이 없으면 뒤의 모든
      숫자가 ``무엇의 평균''인지 모른다. 학습 곡선이 \emph{에피소드
      단위}인지 \emph{스텝 단위} 평균인지도 여기서.
    \item \kb{(2) 알고리즘 선택} — ``PPO를 골랐다''가 아니라,
      \textbf{이 환경의 어떤 특성}이 그 선택을 정당화하는지.
      \begin{itemize}
        \item \textbf{연속 vs 이산} 행동 공간이 \textbf{가장 먼저} 봐야
          하는 분기: DQN(Ch09)은 \emph{이산} 행동을 전제로 $Q(s,a)$를
          표처럼 출력 — 연속(Pendulum, Reacher)에 그대로 쓰면 ``행동을
          어떻게 이산화했는가''의 추가 전제가 붙는다.
        \item \textbf{희소} 보상(Ch15.2의 보너스 구조)도 근거: 희소 보상에
          DQN은 Q값의 \emph{희소성 문제}(대부분 스텝이 0이라 학습 신호
          약함)를 겪는다.
      \end{itemize}
  \end{itemize}
\end{frame}
